\documentclass[12pt,a4paper]{article}

\usepackage{amsmath, amsthm, amssymb, amsfonts}
\usepackage{mathrsfs}
\usepackage{hyperref}
\usepackage{geometry}
\usepackage{enumitem}
\usepackage{booktabs}

\geometry{margin=2.5cm}

\newtheorem{theorem}{Theorem}[section]
\newtheorem{proposition}[theorem]{Proposition}
\newtheorem{lemma}[theorem]{Lemma}
\newtheorem{corollary}[theorem]{Corollary}
\newtheorem{definition}[theorem]{Definition}
\newtheorem{remark}[theorem]{Remark}
\newtheorem{assumption}[theorem]{Assumption}
\newtheorem{openproblem}[theorem]{Open Problem}
\newtheorem{conjecture}[theorem]{Conjecture}

\newcommand{\Ai}{\operatorname{Ai}}
\newcommand{\lam}[1]{\lambda_{#1}^{(c)}}
\newcommand{\eps}[1]{\varepsilon_{#1}^{(c)}}
\newcommand{\pswf}[1]{\psi_{#1}^{(c)}}
\newcommand{\norm}[1]{\left\|#1\right\|}
\newcommand{\Tc}{T_c}
\newcommand{\VN}{V_N}
\newcommand{\Lop}{\mathcal{L}}

\title{%
  \textbf{Paper~IX: B-strong as a Spectral Gap}\\[0.4em]
  \large Operator-Theoretic Reformulation
  in the Superconcentration Regime
}
\author{\textit{prolate-gram-coercivity programme}}
\date{May 2026}

\begin{document}
\maketitle

\begin{abstract}
Assumption B-strong ($P_{kl} \leq C_2 c^{1/2}$, Paper~VI) is
reformulated as a \emph{spectral gap condition} for the
prolate operator $\Tc$ restricted to the Galerkin subspace
$\VN = \operatorname{span}\{\pswf{0},\ldots,\pswf{N}\}$,
$N = \lfloor Ac^{1/3}\rfloor$.
The Airy-WKB analysis of \texttt{bstrong\_reduction\_lemma.tex}
established that the Galerkin indices $k \leq N$ lie entirely
in the \emph{superconcentration regime}
$\varepsilon_k := 1 - \lam{k} \ll 1$,
asymptotically disjoint from the spectral-edge Airy window.
This disjointness, together with the structure of $P_{kl}$
as a ratio of eigenvalue residuals, reduces B-strong to the
statement that the spectral gap
$\gamma(c) := 1 - \|\Tc|_{\VN}\|$
satisfies
\[
  \gamma(c) \geq C c^{-1/2}.
\]
We formulate this as the central open problem of the
programme, identify the natural proof routes
(Landau--Pollak concentration, resolvent methods, direct
coercivity), and explain why the spectral-gap formulation
bypasses the failure modes of the WKB approach.
\end{abstract}

\tableofcontents
\bigskip\hrule\bigskip

%% ============================================================
\section{From Eigenvalue Ratios to Spectral Gap}
\label{sec:reduction}
%% ============================================================

\subsection{The superconcentration regime}

The Galerkin truncation is at $N = \lfloor Ac^{1/3}\rfloor$.
The Slepian concentration number is $N_{\mathrm{Sh}} = 2c/\pi$.
Since $c^{1/3} \ll c$ as $c \to \infty$:
\[
  \frac{N}{N_{\mathrm{Sh}}} \sim \frac{A\pi}{2}\,c^{-2/3} \to 0.
\]
All Galerkin indices $k \leq N$ therefore satisfy
$k \ll N_{\mathrm{Sh}}$, which is the \emph{superconcentration
regime} for PSWF eigenvalues: the concentration eigenvalues
$\lam{k}$ are exponentially close to $1$.

\begin{definition}[Eigenvalue residuals]
For $k \leq N$, define the \emph{residual}
\[
  \eps{k} := 1 - \lam{k} > 0.
\]
By the superconcentration property, $\eps{k} \ll 1$
uniformly for $k \leq N$ and all sufficiently large $c$.
\end{definition}

\subsection{B-strong as a residual ratio condition}

In terms of residuals, the B-strong prefactor is:
\[
  P_{kl} = c \cdot \frac{|\lam{l} - \lam{k}|}{(1-\lam{k})(1-\lam{l})}
  = c \cdot \frac{|\eps{k} - \eps{l}|}{\eps{k}\,\eps{l}}.
\]

\begin{proposition}[B-strong in residual form]
\label{prop:residual}
Assumption B-strong ($P_{kl} \leq C_2 c^{1/2}$ for all
$k,l \leq N$, $|k-l| \leq c^{1/6}$) is equivalent to:
\[
  \frac{|\eps{k} - \eps{l}|}{\eps{k}\,\eps{l}} \leq C_2\,c^{-1/2}
\]
for all such $k, l$.
Alternatively, with $d = |k - l|$ and $k < l$:
\[
  \frac{1}{\eps{k}} - \frac{1}{\eps{l}} \leq C_2\,c^{-1/2},
\]
which is a \emph{Lipschitz condition on the map
$k \mapsto 1/\eps{k}$} with Lipschitz constant $C_2 c^{-1/2}$
(per unit index step, weighted by $d$).
\end{proposition}

\begin{proof}
Direct computation:
$|\eps{k}-\eps{l}|/(\eps{k}\eps{l}) = |1/\eps{l} - 1/\eps{k}|$.
\end{proof}

\begin{remark}[Structural parallel with Airy analysis]
The map $k \mapsto 1/\eps{k}$ is the discrete analogue
of $x \mapsto 1/S(x)$ from \texttt{bstrong\_reduction\_lemma.tex}.
The Lipschitz condition here plays the role of
$\delta \cdot \Phi(x) \leq Cc^{-1/6}$ there.
The difference is that in the superconcentration regime
$\eps{k}$ is not approximated by any simple Airy expression,
so the Lipschitz bound must come from operator-theoretic
or spectral arguments.
\end{remark}

%% ============================================================
\section{Operator-Theoretic Reformulation}
\label{sec:operator}
%% ============================================================

\subsection{The restricted operator and its spectrum}

Let $\Tc : L^2(\mathbb{R}) \to L^2(\mathbb{R})$ be the
prolate operator with kernel
$K_c(x,y) = \sin(c(x-y))/(\pi(x-y))$,
restricted to $[-1,1]$.
The PSWFs $\pswf{k}$ form a complete orthonormal system
of eigenfunctions:
\[
  \Tc \pswf{k} = \lam{k}\,\pswf{k},
  \qquad 1 > \lam{0} > \lam{1} > \cdots > 0.
\]
Define the Galerkin subspace
$\VN := \operatorname{span}\{\pswf{0},\ldots,\pswf{N}\}$
and the restricted operator $\Tc|_{\VN} : \VN \to \VN$.
Since $\VN$ is invariant under $\Tc$ (it is spanned by
eigenvectors), the operator norm is:
\[
  \|\Tc|_{\VN}\| = \lam{0}^{(c)} < 1.
\]

\subsection{The spectral gap}

\begin{definition}[Spectral gap]
The \emph{spectral gap} of $\Tc$ relative to $\VN$ is:
\[
  \gamma(c) := 1 - \|\Tc|_{\VN}\| = 1 - \lam{0}^{(c)} = \eps{0}.
\]
\end{definition}

\begin{remark}
Since $\lam{0} > \lam{1} > \cdots > \lam{N}$,
the relevant spectral gap is controlled by $\eps{0}$,
the residual of the \emph{largest} Galerkin eigenvalue.
B-strong in its strongest form (for all $k, l$) requires
control of the entire ratio $1/\eps{k} - 1/\eps{l}$,
not just $\eps{0}$.
However, the dominant contribution to $P_{kl}$ for
$k \approx 0$, $l \approx N$ comes from $1/\eps{N}$,
the inverse residual of the \emph{smallest} Galerkin eigenvalue.
\end{remark}

\subsection{B-strong as a two-sided spectral condition}

\begin{proposition}[B-strong via spectral gap]
\label{prop:spectral_gap}
A sufficient condition for B-strong is:
\[
  \frac{1}{\eps{N}} - \frac{1}{\eps{0}} \leq C_2\,c^{-1/2}.
\]
This is implied by:
\[
  \frac{d}{dk}\left(\frac{1}{\eps{k}}\right) \leq \frac{C_2\,c^{-1/2}}{d}
  \quad \text{for all } k \leq N,\; d = |k-l|,
\]
i.e.\ $k \mapsto 1/\eps{k}$ is Lipschitz-$c^{-1/2}$ on
$\{0, \ldots, N\}$.
\end{proposition}

\begin{remark}[Why this is non-trivial]
For $k$ in the superconcentration regime,
$\eps{k}$ is exponentially small and the ratio
$1/\eps{k}$ is exponentially large.
B-strong requires that despite this,
the \emph{differences} $1/\eps{l} - 1/\eps{k}$
remain polynomially bounded ($\leq C c^{-1/2}$).
This is a statement of \emph{relative stability}:
the exponentially small residuals vary slowly
\emph{relative to each other}.
\end{remark}

%% ============================================================
\section{The Central Open Problem}
\label{sec:open}
%% ============================================================

\begin{openproblem}[Spectral gap lower bound]
\label{op:gap}
Show that
\[
  \frac{1}{\eps{N}^{(c)}} - \frac{1}{\eps{0}^{(c)}}
  \leq C c^{-1/2}
\]
for all $c \geq c_0$, with $N = \lfloor Ac^{1/3}\rfloor$
and $C$ depending only on $A$.

\medskip
\noindent
Equivalently: the map $k \mapsto 1/\eps{k}^{(c)}$
is Lipschitz-$(C c^{-1/2})$ on $\{0, \ldots, N\}$,
uniformly in $c$.
\end{openproblem}

\subsection{Proof routes}

\subsubsection*{Route A: Landau--Pollak concentration bounds}

Landau and Pollak \cite{LandauPollak1961} established that
for $k \leq N_{\mathrm{Sh}} = 2c/\pi$:
\[
  \lam{k} \geq 1 - \exp(-c\,\phi(k/c))
\]
for an explicit function $\phi > 0$.
For $k \leq N = O(c^{1/3})$:
\[
  \frac{k}{c} \leq Ac^{-2/3} \to 0,
  \qquad \phi(k/c) \sim \phi(0) + O(c^{-2/3}).
\]
A precise asymptotic expansion of $\eps{k}$ in powers
of $k/c$ for $k \ll N_{\mathrm{Sh}}$ would give
explicit Lipschitz bounds for $1/\eps{k}$.

\textbf{Obstacle:} The Landau--Pollak bound is not sharp
enough at this scale; it gives
$\eps{k} \sim \exp(-c\phi(k/c))$
but not the precise $k$-dependence of the exponent
at $k = O(c^{1/3})$.

\subsubsection*{Route B: Eigenvalue perturbation / resolvent}

Since $\lam{k}$ are simple eigenvalues of $\Tc$,
the Hadamard--Weyl formula gives:
\[
  \lam{k} - \lam{k+1}
  = \langle \Tc\pswf{k}, \pswf{k}\rangle
  - \langle \Tc\pswf{k+1},\pswf{k+1}\rangle
  = \int_{-1}^{1}\int_{-1}^{1}
    K_c(x,y)(\pswf{k}(x)\pswf{k}(y)
    - \pswf{k+1}(x)\pswf{k+1}(y))\,dx\,dy.
\]
For $k \leq N$ in the superconcentration regime,
$\pswf{k}$ and $\pswf{k+1}$ are both near-identical
functions (both highly concentrated in $[-1,1]$),
and their difference controls $\lam{k}-\lam{k+1}$.
A resolvent estimate on $\Tc - \lam{k}$ for
$k$ in the superconcentration regime would quantify this.

\textbf{Obstacle:} The resolvent of $\Tc$ is not
explicit; one needs spectral theory of
integral operators with oscillatory kernels.

\subsubsection*{Route C: Direct coercivity on $\VN$}

Instead of bounding individual $\eps{k}$, one can
attempt to show directly that the bilinear form
associated to the Galerkin matrix $G_N$ (Paper~VI)
is coercive with constant $\geq C c^{-1/2}$
on $\VN$, without passing through $P_{kl}$.
This bypasses B-strong entirely and goes directly
to the target estimate of the programme.

Specifically: if one can show
\[
  \langle (I - \Tc)u, u\rangle \geq C c^{-1/2}\,\|u\|^2
  \quad \forall\, u \in \VN,
\]
then $\gamma(c) \geq Cc^{-1/2}$ follows immediately.
For $u = \sum_{k=0}^N a_k \pswf{k}$:
\[
  \langle (I-\Tc)u,u\rangle
  = \sum_{k=0}^N \eps{k}\,|a_k|^2
  \geq \eps{N}\,\|u\|^2.
\]
Hence the coercivity constant is exactly $\eps{N}$,
and B-strong (in its strongest form) is equivalent to
$\eps{N} \geq C c^{-1/2}$ \emph{and} to the
Lipschitz condition on $1/\eps{k}$.

\begin{proposition}[Coercivity equivalent]
\label{prop:coercivity}
$(I - \Tc)|_{\VN}$ is coercive with constant
$\eps{N} = 1 - \lam{N}^{(c)}$.
B-strong (for all $k, l \leq N$, $d = |k-l| \leq c^{1/6}$)
is implied by:
\[
  \eps{N}^{(c)} \geq c^{1/2} \cdot C^{-1}
  \qquad \text{(}\!\Leftarrow\!\text{ coercivity at scale } c^{-1/2}\text{).}
\]

\noindent\emph{Warning:} this is a one-sided statement.
B-strong also requires \emph{upper} bounds on
$\eps{k}$ differences (the Lipschitz condition),
not only lower bounds on $\eps{N}$.
See Open Problem~\ref{op:gap}.
\end{proposition}

%% ============================================================
\section{Known Asymptotics of $\varepsilon_k$ and What Is Missing}
\label{sec:asymptotics}
%% ============================================================

\subsection{Qualitative picture}

The following is known for PSWF eigenvalues
(see Osipov--Rokhlin--Xiao \cite{ORX2013},
Slepian \cite{Slepian1978}):
\begin{enumerate}[nosep]
  \item For $k \leq N_{\mathrm{Sh}} - O(c^{1/3})$
        (well below the spectral edge):
        $\eps{k} = 1 - \lam{k}$ is superexponentially
        small in $c$.
  \item Near the edge ($k \approx N_{\mathrm{Sh}}$):
        $\eps{k}$ transitions from $\approx 0$ to
        $\approx 1$ on a scale $O(c^{1/3})$,
        described by the Airy CDF $F_{\mathrm{Ai}}$.
  \item For $k \leq N = O(c^{1/3})$:
        $\eps{k} \sim \exp(-c^\alpha g(k/c))$
        for some $\alpha > 0$ and $g > 0$;
        the precise exponent is not pinned down
        in the literature to sufficient precision for B-strong.
\end{enumerate}

\subsection{What is missing}

\begin{openproblem}[Superconcentration asymptotics]
\label{op:asymptotics}
Determine the asymptotic behaviour of
$\eps{k}^{(c)} = 1 - \lam{k}^{(c)}$
for $k \leq N = \lfloor Ac^{1/3}\rfloor$
as $c \to \infty$, to sufficient precision to prove
or disprove the Lipschitz condition of
Open Problem~\ref{op:gap}.

\medskip\noindent
Specifically: is there an expansion
\[
  \eps{k}^{(c)} = \exp\!\bigl(-c^{\alpha}\,h(k, c)\bigr)
\]
where $h(k,c)$ varies slowly enough in $k$ that
$1/\eps{k} - 1/\eps{l} \leq Cc^{-1/2}$
for $|k-l| \leq c^{1/6}$?
\end{openproblem}

%% ============================================================
\section{Structural Map: Old vs.\ New Framework}
\label{sec:map}
%% ============================================================

\begin{center}
\renewcommand{\arraystretch}{1.5}
\begin{tabular}{p{3.5cm} p{4.5cm} p{4.5cm}}
\toprule
& \textbf{Airy-WKB framework}
& \textbf{Operator / superconc.\ framework} \\
\midrule
Primary object
& $S(x) = \int_x^\infty \Ai^2$, $\Phi(x) = \Ai^2/S^2$
& $\eps{k} = 1 - \lam{k}$, map $k \mapsto 1/\eps{k}$ \\
Key identity
& $d(1/S)/dx = \Phi(x)$
& $|1/\eps{l} - 1/\eps{k}| = P_{kl}/c$ \\
Target bound
& $\delta\,\Phi(x) \leq Cc^{-1/6}$
& $1/\eps{N} - 1/\eps{0} \leq Cc^{-1/2}$ \\
Failure mode
& $\Phi$ not uniform in transition regime
& $1/\eps{k}$ exponentially large (but differences small?) \\
Relevant spectral domain
& $x \in [-A_0, (\log c)^{2/3}]$ ($k \sim N_{\mathrm{Sh}}$)
& $k \leq N = O(c^{1/3})$ ($x_k \sim -c^{2/3}$) \\
Domains disjoint?
& \multicolumn{2}{c}{\textbf{Yes.} $N/N_{\mathrm{Sh}} \sim c^{-2/3} \to 0$.} \\
Tool
& Airy ODE, WKB, Wronskian
& Eigenvalue asymptotics, resolvent, coercivity \\
Status
& Closed in own domain; spectrally empty for B-strong
& \textbf{Open.} Central problem of this note. \\
\bottomrule
\end{tabular}
\end{center}

%% ============================================================
\section{Relation to the Programme and Next Steps}
\label{sec:next}
%% ============================================================

This paper represents a \emph{direction shift} in the
prolate-gram-coercivity programme.
Papers~I--VIII developed analytical tools
(Airy approximation, WKB, semiclassical phase,
scale-separated estimates) whose natural domain
is the spectral-edge regime $k \sim N_{\mathrm{Sh}}$.
Paper~IX identifies that B-strong, for the Galerkin
indices $k \leq N = O(c^{1/3})$, requires instead
tools from the spectral theory of
$\Tc|_{\VN}$ in the superconcentration regime.

\paragraph{Immediate next step.}
The most direct route to B-strong is
Route~C (Section~\ref{sec:open}):
show $\eps{N}^{(c)} \geq Cc^{-1/2}$ using
the known concentration properties of $\pswf{N}$
for $N = O(c^{1/3})$.
This is a sharp version of the Landau--Pollak theorem
at the Galerkin scale, and is likely accessible
via the integral-equation methods of
Osipov--Rokhlin--Xiao \cite{ORX2013}, Chapter~4.

\paragraph{Alternative.}
If $\eps{N} \geq Cc^{-1/2}$ proves elusive,
one may attempt to bypass B-strong entirely
(Route~C coercivity without the $P_{kl}$ decomposition)
or to establish a weaker uniform bound
$P_{kl} \leq Cc^\alpha$ for some $\alpha < 1$
that still suffices for the coercivity conclusion
of Paper~VI.

%% ============================================================
\begin{thebibliography}{9}

\bibitem{paper6}
  \textit{Paper~VI: Galerkin Norm Estimates for the Prolate Phase Operator},
  \texttt{paper6.tex}, prolate-gram-coercivity programme, May~2026.

\bibitem{bstrong_reduction}
  \textit{B-strong Reduction Lemma},
  \texttt{bstrong\_reduction\_lemma.tex},
  prolate-gram-coercivity programme, May~2026.

\bibitem{LandauPollak1961}
  H.~J.~Landau and H.~O.~Pollak,
  ``Prolate Spheroidal Wave Functions, Fourier Analysis and
  Uncertainty~II,''
  \textit{Bell Syst.\ Tech.\ J.}, 40 (1961), 65--84.

\bibitem{Slepian1978}
  D.~Slepian,
  ``Prolate Spheroidal Wave Functions, Fourier Analysis, and
  Uncertainty~V,''
  \textit{Bell Syst.\ Tech.\ J.}, 57 (1978), 1371--1430.

\bibitem{ORX2013}
  A.~Osipov, V.~Rokhlin, H.~Xiao,
  \textit{Prolate Spheroidal Wave Functions of Order Zero},
  Springer, 2013.

\end{thebibliography}

\end{document}
